\subsubsection{The extension to arbitrary PL manifolds}

We retain the convention that a concordance of a compact manifold $M$
fixes $M\times\{0\}\cup\partial M\times I$ pointwise, with the top face
free.  PL components are defined by jointly PL families with jointly PL
inverses, and their group law is composition.  For a compact manifold
$X$, write $E(X)$ for the previously defined countable presentation
formed from a finite homotopy model of $X$.

\begin{lemma}\label{lem:component-followup-E-naturality}
If $a:X\to Y$ is a based three-connected map of connected finite CW
homotopy types, the induced homomorphism
\[
                         E(a):E(X)\longrightarrow E(Y)
\]
is an isomorphism.  In particular, $E$ is invariant under homotopy
equivalence and $E(K^{(3)})\to E(K)$ is an isomorphism for every connected
finite simplicial complex $K$.
\end{lemma}

\begin{proof}
The loop-group map is two-connected.  Its map of connective spherical
group rings is therefore two-connected.  At each finite rank, extension
of scalars induces a two-connected map of derived endomorphism spaces
and an isomorphism of their degree-zero matrix rings.  It identifies
the invertible components, so the homotopy-automorphism monoids receive
a two-connected map and their bar constructions receive a
three-connected map.  The same conclusion holds for the bars of the
monomial groups: on each permutation component, the underlying map is
a finite product of the loop-group map.

These maps commute with adding an identity summand and with the
monomial inclusion.  The corresponding maps of the two telescopes are
three-connected.  They identify fundamental groups and their
abelianization kernels, and hence lift to three-connected maps of the
specified abelian covers.  In particular, the maps on integral second
homology are isomorphisms.  The presentation $E$ is the cokernel of the
monomial map on these second homology groups, so its induced map is an
isomorphism.  The argument uses no simple-connectivity assumption on
the abelian covers.

Homotopy equivalences are three-connected.  The inclusion
$K^{(3)}\to K$ is three-connected because its omitted cells have
dimensions at least four.  This proves the last assertions.
\end{proof}

\begin{lemma}\label{lem:component-followup-handle-core}
Let $M$ be a compact connected PL manifold of dimension $n\geq6$.
There is a smoothable compact codimension-zero PL submanifold
$j:N\hookrightarrow\operatorname{int}M$ having the homotopy type of a
finite CW complex of dimension at most three, such that
\[
                         \pi_i(M,N)=0\quad(i\leq3),
\]
and $\partial N\to N$ induces isomorphisms on $\pi_0$ and $\pi_1$.
The manifold $N$ is connected.
\end{lemma}

\begin{proof}
Take an absolute finite handle decomposition of $M$, relative to the
empty set, with indices ordered increasingly, and let $N_0$ be the
stage through index three.  The triangulation construction supplies
such a decomposition, with one handle of index $\dim A$ for each
simplex $A$; its associated core complex has one cell for each handle
\cite[Proposition~6.9 and the following discussion,
pp.~82--83]{RourkeSanderson1982Component}.

If $M$ has boundary, push a compact PL copy of $M$ into its interior
using a boundary collar and let $N$ be the corresponding copy of
$N_0$.  This collar embedding is homotopic to the identity as a map to
$M$.  For closed $M$ take $N=N_0$.  Adding the usual collars to handle
stages does not change their indices or their homotopy types.
Thus $N$ has the homotopy type of a finite three-complex.  The remaining
handles have indices at least four, which proves
$\pi_i(M,N)=0$ for $i\leq3$.  They cannot join components, so $N$ is
connected.

The dual handle decomposition of $N$ relative to its entire boundary
has indices at least $n-3\geq3$; see the duality construction on
page~82 of the same source.  These attachments preserve components
and fundamental groups.  Hence $\partial N\to N$ is an isomorphism on
$\pi_0,\pi_1$.

Finally, the fiber $PL/O$ of the stable smoothing lift is
6-connected.  All lifting obstructions vanish for a PL manifold with
the homotopy type of a three-dimensional CW complex.  The smoothing
theorem therefore gives a compatible smoothing
\cite[Part~II, Theorem~4.1, p.~108, and \S5.3,
pp.~122--123]{HirschMazur1974}.
For boundary, apply the lifting theorem to the interior, then use a
collar to take a compact PL copy inside that interior and move its
locally flat boundary to a smooth hypersurface
\cite[Part~I, Theorem~7.4, pp.~37--38]{HirschMazur1974}.
The resulting compact domain gives a compatible smoothing of $N$.
No smoothing of $M$ or preassigned smoothing of $\partial N$ is required.
\end{proof}

\begin{lemma}\label{lem:component-followup-stable-square}
For a codimension-zero inclusion $j:N\hookrightarrow\operatorname{int}M$,
extension by the identity induces a homomorphism
$\gamma_0:\pi_0P_{\mathrm{PL}}(N)\to\pi_0P_{\mathrm{PL}}(M)$.
If $s_X:\pi_0P_{\mathrm{PL}}(X)\to E(X)$ denotes stabilization followed
by the established stable comparison, then
\[
\begin{array}{ccc}
\pi_0P_{\mathrm{PL}}(N)&\xrightarrow{\ \gamma_0\ }&
       \pi_0P_{\mathrm{PL}}(M)\\
\downarrow s_N&&\downarrow s_M\\
E(N)&\xrightarrow{\ E(j)\ }&E(M)
\end{array}
\tag{$*$}\label{eq:component-followup-stable-square}
\]
commutes.
\end{lemma}

\begin{proof}
On a parameter simplex, extend the given map on
$\Delta^q\times N\times I$ by the identity.  The two definitions agree
on $\Delta^q\times\partial N\times I$.  Common finite subdivisions
therefore give a jointly PL map and its jointly PL inverse on the
whole product.  The parameter projection, bottom fixation, and side
fixation are retained.  Extension commutes with composition and
inversion, which proves the first assertion.

For naturality, use the collared h-cobordism model of the stable
comparison.  A codimension-zero embedding induces the covariant
pushout
\[
                    W\longmapsto M\times I\cup_{N\times I}W
\]
on collared h-cobordisms
\cite[Definition~1.1.3(b), pp.~8--9]{WaldhausenJahrenRognes2013}.
On automorphisms of the product h-cobordism, this is exactly extension
by the identity.  The product component is the classifying space of
the concordance group with the bottom-and-side convention used here;
thus its fundamental group identifies with the concordance component
group with composition
\cite[Remark~1.1.11 and (1.1.12),
pp.~12--13]{WaldhausenJahrenRognes2013}.

The stable h-cobordism comparison is natural for these maps
\cite[Theorem~0.1, p.~1, and Theorem~1.1.8,
p.~11]{WaldhausenJahrenRognes2013}.
In the retractive-space model its map is pushout along $j$; in the
spherical-group-ring model it is extension of scalars.  This preserves
the monomial maps and rank stabilization and induces the maps of the
based abelian covers and their second homology cokernels.  Consequently
the induced map on the specified presentation is $E(j)$.
For the PL comparison, the assembly cofiber sequence in the proof of
Theorem~\ref{thm:stable-pl-component-presentation} is natural in the
base. Its isomorphism from the degree-two smooth Whitehead group to
the degree-two PL Whitehead group, and therefore its inverse, commute
with these maps. This verifies the same identification for a possibly
nonsmoothable $M$; no arbitrary choice of a group isomorphism is used.

Positive suspension commutes with extension by the identity on
components: choose collar-normalized representatives and the same
positive bending in the interval coordinates.  Both composites agree
on the $N$-part and are the identity outside it.  The same argument
applies to compact families.  Passing to stabilization colimits and
then to the natural stable comparison gives the displayed square.
\end{proof}

\begin{theorem}\label{thm:component-followup-all-PL-six}
For every compact connected PL manifold $M$ of dimension at least six,
stabilization induces an isomorphism
\[
                      \pi_0P_{\mathrm{PL}}(M)\xrightarrow{\ \cong\ }E(M).
\]
Consequently, for every connected finite simplicial complex $K\simeq M$,
\[
                             \pi_0P_{\mathrm{PL}}(M)\cong E(K).
\]
No compatible smoothing of $M$ is assumed.
\end{theorem}

\begin{proof}
Choose $N$ by Lemma~\ref{lem:component-followup-handle-core}.
We need only that extension by the identity is surjective on
components.  The original component reduction to an absolute
three-handle stage is explicitly recorded in Pedersen's appendix
\cite[\S2, pp.~148--149]{Pedersen1975ComponentAppendix}.
For the precise extension map and a direct check of the required
weak consequence, use the PL statement in
\cite[Theorem~3.3 and the preceding paragraph,
p.~17]{GalvinNonino2025v1Component}.
That statement assumes
\[
 \begin{gathered}
 \pi_i(\partial N)\cong\pi_i(N)\ (i=0,1),\qquad
 \pi_i(M,N)=0\ (i\leq r),\\
 N\text{ is }k\text{-connected},\qquad 0\leq k\leq r\leq n-4.
 \end{gathered}
\]
Its PL connectivity bound for the extension map is
\[
                              \min(2r-3,r+k-2,r+2).
\]
Take $r=2$ and $k=0$.  All the hypotheses hold, and the bound is
$\min(1,0,4)=0$, since $2\leq n-4$.  Hence $\gamma_0$ is surjective.
This use requires no hypothesis on $\partial M\to M$ and does not use
an $r=3$ application in dimension six.

The smoothable dimension-six result established above applies to
$N$ and gives that $s_N$ is an isomorphism.
Lemma~\ref{lem:component-followup-E-naturality} gives that $E(j)$ is
an isomorphism because $j$ is three-connected.
The natural square in Lemma~\ref{lem:component-followup-stable-square}
therefore gives
\[
                              s_M\gamma_0=E(j)s_N
\]
as an isomorphism.  If $\gamma_0(a)=\gamma_0(b)$, applying $s_M$
and the injectivity of the right-hand side gives $a=b$.
Thus $\gamma_0$ is also injective, and $s_M$ is an isomorphism.
Homotopy invariance of $E$ proves the assertion for $K$.

\end{proof}

\begin{corollary}\label{cor:component-followup-chosen-three-skeleton}
Let $M$ be a compact connected PL manifold of dimension at least six,
possibly with boundary, and let $K$ be a specified finite
combinatorial triangulation of $M$.  Then
\[
                              \pi_0P_{\mathrm{PL}}(M)\cong E(K^{(3)}).
\]
\end{corollary}

\begin{proof}
The theorem gives $\pi_0P_{\mathrm{PL}}(M)\cong E(K)$.
The inclusion $K^{(3)}\to K$ is three-connected, so
Lemma~\ref{lem:component-followup-E-naturality} gives
$E(K^{(3)})\cong E(K)$.
This proves the conclusion for the specified skeleton.  It does not
require identifying the particular handle core used in the proof of
the theorem with that skeleton. 
\end{proof}

\begin{corollary}\label{cor:component-followup-fields-products}
The same presentations compute the stipulated definable concordance
component groups for finite PL models over every real closed field
and in each fixed o-minimal expansion.  For a compact manifold with
finitely many components, use the finite product of the connected
presentations, with each component in the corresponding dimension
range.
\end{corollary}

\begin{proof}
Use the established relative finite-model PL/definable component
comparison and its field independence, which preserve composition
and the entire fixed bottom-and-side subpolyhedron.
Bottom fixation forces each concordance and each permitted family
to preserve every connected component of the manifold.  Consequently
the groups and their component groups decompose as the asserted
finite products.
\end{proof}
